\documentclass[../../../../main.tex]{subfiles}
\begin{document}

\begin{flushright}
\footnotesize\textit{【自動文字起こし・要確認】}
\end{flushright}

[解] $f(x) = k(x^{m+1}-1) - (x^m-1)$ とおく。$f(x) \ge 0$ が $0 < x$ の任意の$x$で成り立つ$k$の条件をしらべる
\begin{align*}
f'(x) &= k(m+1)x^m - mx^{m-1} \\
&= \{k(m+1)x - m\} x^{m-1}
\end{align*}
から、下表を得る

\begin{center}
\begin{tabular}{|c|c|c|c|c|}
\hline
$x$ & $0$ & $\cdots$ & $\frac{m}{k(m+1)}$ & $\cdots$ \\ \hline
$f'$ & & $-$ & $0$ & $+$ \\ \hline
$f$ & & $\searrow$ & & $\nearrow$ \\ \hline
\end{tabular}
\end{center}

よって、
\begin{align*}
f(x) &\ge f\left(\frac{m}{k(m+1)}\right) = k \left[ \left\{\frac{m}{k(m+1)}\right\}^{m+1} - 1 \right] - \left[ \left\{\frac{m}{k(m+1)}\right\}^m - 1 \right] \\
&= \left\{\frac{m}{k(m+1)}\right\}^m \left( \frac{m}{m+1} - 1 \right) + (1-k) \cdots \text{①}
\end{align*}

$x \to \infty$ とすれば、$k > 0$ が必要なので、$t = \frac{1}{k}$ とおいて、①を$g(t)$とすると
\[ g(t) = - \frac{m^m}{(m+1)^{m+1}} t^m + 1 - \frac{1}{t} \]
\[ g'(t) = - \frac{m^{m+1}}{(m+1)^{m+1}} t^{m-1} + \frac{1}{t^2} \]
より、下表をえる

\begin{center}
\begin{tabular}{|c|c|c|c|c|}
\hline
$t$ & $0$ & $\cdots$ & $\frac{m+1}{m}$ & $\cdots$ \\ \hline
$g'$ & & $+$ & $0$ & $-$ \\ \hline
$g$ & & $\nearrow$ & & $\searrow$ \\ \hline
\end{tabular}
\end{center}

これから
\[ g(t) \le g\left(\frac{m+1}{m}\right) = 0 \]
だから、常に $f(x) \ge 0$ がみたされるには、$t = \frac{m+1}{m} \Leftrightarrow k = \frac{m}{m+1}$ が条件である。

\end{document}
